\chapter{Hardware and Connections}\label{ch:hardware}

\section{Parts}

Table~\ref{tab:parts} lists the parts of a ClearHPGL plotter. The mechanics
(the frame, the guides, and the belts or screws) are your design. The
firmware drives two axes and a pen lift, so the plot area is the travel that
your mechanics give.

\begin{table}[htbp]
\centering
\caption{Parts of a ClearHPGL plotter}\label{tab:parts}
\begin{tabularx}{\linewidth}{@{}P{0.3\linewidth}L@{}}
\toprule
\textbf{Part} & \textbf{Function and notes} \\
\midrule
Teknic ClearCore & Runs the firmware. Needs its own 24~V DC supply. \\
Two ClearPath-SD servos & Step-and-direction servos for X and Y. Each one
needs its HLFB (ready) output connected. \\
Pen lift & A solenoid or a relay on a digital output, or a third ClearPath
servo on \conn{M2}. \\
Emergency stop & A normally closed contact on \conn{DI-6}, and a separate
contact that removes power from the motors. \\
Network cable or USB cable & The link to the computer. A serial cable on
\conn{COM-0} or \conn{COM-1} is also possible. \\
Windows computer & For the desktop app and the CAD software. The app is
optional. \\
\bottomrule
\end{tabularx}
\end{table}

\section{Connector Map}

Table~\ref{tab:connectors} shows what each connector of the ClearCore does
with the default settings.

\begin{table}[htbp]
\centering
\caption{ClearCore connectors used by ClearHPGL}\label{tab:connectors}
\begin{tabularx}{\linewidth}{@{}P{0.16\linewidth}P{0.28\linewidth}L@{}}
\toprule
\textbf{Connector} & \textbf{Function} & \textbf{Setting} \\
\midrule
\conn{M0} & X axis (CoreXY: motor A), ClearPath, step and direction & \key{steps_rev_x}, \key{mm_rev_x}, \key{dir_x}, \key{kinematics} \\
\conn{M1} & Y axis (CoreXY: motor B), ClearPath, step and direction & \key{steps_rev_y}, \key{mm_rev_y}, \key{dir_y}, \key{swap_xy} \\
\conn{M2} & Pen lift servo (optional) & \key{pen_mode}~1 and the \key{z_*} settings \\
\conn{IO-0} & Pen lift output (default) & \key{pen_mode}~0, \key{pen_io}, \key{pen_active_high} \\
\conn{DI-6} & Emergency stop input & \key{estop_di6} \\
\conn{COM-0}, \conn{COM-1} & HP-GL serial port (off by default) & \key{com_port}, \key{com_baud}, \key{com_handshake}, \key{com_rs232} \\
USB & HP-GL and the setup console & none \\
Ethernet & HP-GL (TCP 9100), console (TCP 9300), discovery (UDP 9302) & \key{net_mode}, \key{ip}, \key{netmask}, \key{gateway} \\
\bottomrule
\end{tabularx}
\end{table}

\begin{figure}[htbp]
\centering
\begin{tikzpicture}[font=\sffamily\small,
  cc/.style={draw=navy,fill=navy!6,rounded corners=3pt,minimum width=2.8cm,minimum height=8.2cm},
  dev/.style={draw=steel,fill=steel!8,rounded corners=2pt,align=center,text width=3.3cm,minimum height=1.0cm,inner sep=3pt},
  pc/.style={draw=rulegray,fill=codebg,rounded corners=2pt,align=center,text width=2.5cm,minimum height=1.0cm,inner sep=3pt},
  conpin/.style={font=\sffamily\scriptsize\bfseries,color=navy},
  w/.style={thick,color=navy},
  ws/.style={thick,color=signal}]
  \node[cc] (cc) {};
  \node[font=\sffamily\bfseries\color{navy}] at ([yshift=-0.35cm]cc.north) {ClearCore};
  % right side: motion
  \node[conpin,anchor=east] (m0) at ([yshift=2.9cm]cc.east) {M0};
  \node[conpin,anchor=east] (m1) at ([yshift=1.45cm]cc.east) {M1};
  \node[conpin,anchor=east] (m2) at ([yshift=0.0cm]cc.east) {M2};
  \node[conpin,anchor=east] (io0) at ([yshift=-1.45cm]cc.east) {IO-0};
  \node[conpin,anchor=east] (di6) at ([yshift=-2.9cm]cc.east) {DI-6};
  \node[dev,anchor=west] (x) at ([xshift=2.5cm]m0.east) {ClearPath-SD\\\scriptsize X axis};
  \node[dev,anchor=west] (y) at ([xshift=2.5cm]m1.east) {ClearPath-SD\\\scriptsize Y axis};
  \node[dev,anchor=west,dashed] (z) at ([xshift=2.5cm]m2.east) {ClearPath (optional)\\\scriptsize pen lift, pen\_mode 1};
  \node[dev,anchor=west] (pen) at ([xshift=2.5cm]io0.east) {Relay or solenoid\\\scriptsize pen lift, pen\_mode 0};
  \node[dev,anchor=west,draw=signal,fill=signal!6] (es) at ([xshift=2.5cm]di6.east) {Emergency stop\\\scriptsize normally closed};
  \draw[w] (m0.east) -- node[above,font=\sffamily\scriptsize]{step, dir, HLFB} (x.west);
  \draw[w] (m1.east) -- node[above,font=\sffamily\scriptsize]{step, dir, HLFB} (y.west);
  \draw[w,dashed] (m2.east) -- node[above,font=\sffamily\scriptsize]{step, dir} (z.west);
  \draw[w] (io0.east) -- node[above,font=\sffamily\scriptsize]{output} (pen.west);
  \draw[ws] (di6.east) -- node[above,font=\sffamily\scriptsize]{input} (es.west);
  % left side: links
  \node[conpin,anchor=west] (usb) at ([yshift=1.6cm]cc.west) {USB};
  \node[conpin,anchor=west] (eth) at ([yshift=0.0cm]cc.west) {Ethernet};
  \node[conpin,anchor=west] (com) at ([yshift=-1.6cm]cc.west) {COM-0/1};
  \node[pc,anchor=east] (pcu) at ([xshift=-1.0cm]usb.west) {Computer\\\scriptsize USB virtual COM};
  \node[pc,anchor=east] (pce) at ([xshift=-1.0cm]eth.west) {Network\\\scriptsize TCP 9100, 9300};
  \node[pc,anchor=east,dashed] (pcc) at ([xshift=-1.0cm]com.west) {Serial host\\\scriptsize RS-232 or TTL};
  \draw[w] (pcu.east) -- (usb.west);
  \draw[w] (pce.east) -- (eth.west);
  \draw[w,dashed] (pcc.east) -- (com.west);
\end{tikzpicture}
\caption{Connections of a ClearHPGL plotter. Dashed parts are optional.}\label{fig:wiring}
\end{figure}

\section{ClearPath Servos}

Connect the X servo to \conn{M0} and the Y servo to \conn{M1}. Use the
ClearPath MSP software from Teknic to set each servo before you connect it.

\begin{procedure}{Prepare a ClearPath-SD servo}
\begin{enumerate}
  \item Connect the servo to the computer with its USB cable, and start
        ClearPath MSP.
  \item Set the input format to step and direction.
  \item Set the input resolution (the steps per revolution). Use the same
        value later in \key{steps_rev_x} or \key{steps_rev_y}. The default is 800.
  \item Set the HLFB output to a mode with a PWM output, with a carrier
        frequency of 482~Hz. The ClearCore examples from Teknic use
        ``ASG-Position w/Measured Torque.''
  \item Tune the servo for the load of its axis, as the ClearPath manual
        tells you.
  \item Save the configuration to the servo.
\end{enumerate}
\end{procedure}

The firmware reads the HLFB signal of each servo. When the motors enable,
the board waits for HLFB from all servos. If HLFB does not come, the board
disables the motors and reports a fault. Set \key{hlfb_check} to 0 only if
your servos do not supply HLFB.

\begin{note}
\notelabel Use the same steps per millimeter on X and Y. The planner uses the
scale of the X axis for the speed and the acceleration of the path. Arcs
operate as true arcs only when the two scales are the same.
\end{note}

\section{CoreXY and H-Bot Machines}\label{sec:corexy}

In a \term{CoreXY} or an H-bot machine, the two motors do not each drive one
axis. Two belts move the head together: when both motors turn the same
direction, the head moves in X; when they turn in opposite directions, it
moves in Y. One motor alone moves the head diagonally.

Set \key{kinematics} to 1 for these machines. The board then plans and draws
in X and Y as usual, and sends these step counts to the motors:

\begin{center}
\begin{tabular}{@{}lll@{}}
\toprule
\textbf{Connector} & \textbf{Motor} & \textbf{Steps} \\
\midrule
\conn{M0} & A & X + Y \\
\conn{M1} & B & X $-$ Y \\
\bottomrule
\end{tabular}
\end{center}

\begin{itemize}
  \item Both belts run on pulleys with the same pitch diameter. Set it in
        \key{pulley_dia_mm}. The travel of the belt for one revolution is
        $\pi$ multiplied by the pitch diameter; \key{mm_rev_x} and
        \key{mm_rev_y} are not used. The default, \SI{12.7324}{mm}, is a
        20-tooth GT2 pulley (\SI{40}{mm} per revolution).
  \item The pitch diameter is the number of teeth multiplied by the belt
        pitch, divided by $\pi$. For example, a 16-tooth GT2 pulley (2~mm
        pitch) has a pitch diameter of 32/$\pi$ = \SI{10.186}{mm}.
  \item Both motors must have the same steps per revolution. The board
        refuses \key{kinematics}~1 if \key{steps_rev_x} and \key{steps_rev_y}
        are different.
  \item On a 45° move, one motor does all the work. It turns 1.41 times faster
        than the head moves, and accelerates 1.41 times faster. Set
        \key{motor_max_mm_s} to the highest belt speed that your motors can
        hold. The board then decreases the speed of each segment that would
        make a motor go faster.
  \item \key{corner_mm_s} limits the change of speed of each motor at a corner.
\end{itemize}

The direction of each motor and the sequence of the belts change how the
head moves. The settings \key{dir_x}, \key{dir_y}, and \key{swap_xy}
(exchange X and Y) correct all the possible layouts. You do not have to
change the wiring. Section~\ref{sec:corexy-check} gives the procedure.

\section{Pen Lift}

The firmware has two types of pen lift. Select one with \key{pen_mode}.

\subsection{Digital Output (\key{pen_mode}~0)}

The default. A digital output (\conn{IO-0} by default) operates a relay or a
solenoid. With \key{pen_active_high}~1, the output is on when the pen is
down. After the board lowers or raises the pen, it waits for the settle time
(\key{pen_down_ms}, default 150~ms; \key{pen_up_ms}, default 100~ms) before
the axes move.

\begin{caution}
Examine the current rating of the ClearCore output in the ClearCore manual.
If the solenoid needs more current, use a relay. Install a flyback diode
across a coil that has no internal diode.
\end{caution}

\subsection{Servo on M2 (\key{pen_mode}~1)}

A third ClearPath servo on \conn{M2} moves the pen to two positions:
\key{z_up_steps} (default 0) and \key{z_down_steps} (default 800). Its speed
and acceleration are \key{z_vel_steps} and \key{z_accel_steps}.

\begin{caution}
The M2 servo has no home position. Before you supply power, put the pen at
the pen-up position. The board uses the position at power-up as step 0.
\end{caution}

\section{Emergency Stop}

Connect a normally closed emergency stop contact to \conn{DI-6}. With the
default \key{estop_di6}~1, the board sees a fault when the input is low
(the contact is open, or a wire is broken). Table~\ref{tab:estop} shows the
three values.

\begin{table}[htbp]
\centering
\caption{Values of \key{estop_di6}}\label{tab:estop}
\begin{tabularx}{\linewidth}{@{}P{0.12\linewidth}L@{}}
\toprule
\textbf{Value} & \textbf{Result} \\
\midrule
0 & The board does not monitor \conn{DI-6}. Use only if no emergency stop is connected. \\
1 & Fault when \conn{DI-6} is low. Use with a normally closed contact (default). \\
2 & Fault when \conn{DI-6} is high. Use with a normally open contact. \\
\bottomrule
\end{tabularx}
\end{table}

\section{Links to the Computer}

\begin{description}[style=nextline,leftmargin=1.2em,font=\sffamily\bfseries\color{navy}]
  \item[USB] The board shows as a virtual COM port. It receives HP-GL at any
        baud rate. Type \cmd{menu} on a line by itself to open the setup
        console.
  \item[Ethernet] The board uses DHCP by default. If DHCP fails, the address is
        192.168.0.110. It receives HP-GL on TCP port 9100 and gives the
        console on TCP port 9300.
  \item[COM-0 or COM-1] An RS-232 or TTL serial port for HP-GL. It is off by
        default. Set \key{com_port}, \key{com_baud}, \key{com_handshake}, and
        \key{com_rs232}, then restart the board.
\end{description}
